% Canonical Paper II source.
Because the Laplace--Beltrami equation in
\eqref{eq:hyperbolic-laplace-beltrami} differs from the Euclidean Laplace
equation only by the positive factor $\lambda^2y^2$, the ideal-boundary
Dirichlet problem is governed by the ordinary upper-half-plane Poisson kernel.
The content of this section is the exact transport to the AEG normalization
and a domain statement that includes the point at infinity.

\subsection{The kernel in both boundary coordinates}

For $X,t\in\R$ and $y>0$, put
\begin{equation}\label{eq:standard-poisson-kernel-X}
 P_y(X-t)=\frac1\pi\frac{y}{(X-t)^2+y^2}.
\end{equation}
For the original $x$-coordinate, define
\begin{equation}\label{eq:aeg-poisson-kernel-x}
 P^{\mu,\lambda}_y(x-\xi)
 :=\frac1\pi
   \frac{(\mu/\lambda)y}
        {(x-\xi)^2+(\mu y/\lambda)^2}
 =\frac{\mu\lambda y}
        {\pi\bigl(\lambda^2(x-\xi)^2+\mu^2y^2\bigr)}.
\end{equation}

\begin{lemma}[Kernel transport and normalization]
\label{lem:poisson-kernel-transport}
If $X=(\lambda/\mu)x$ and $t=(\lambda/\mu)\xi$, then
\begin{equation}\label{eq:poisson-density-transport}
 P_y(X-t)\,dt
 =P^{\mu,\lambda}_y(x-\xi)\,d\xi.
\end{equation}
In particular,
\[
 P^{\mu,\lambda}_y\geq0,
 \qquad
 \int_{\R}P^{\mu,\lambda}_y(x-\xi)\,d\xi=1.
\]
For fixed $\xi$, the function $(x,y)\mapsto
P^{\mu,\lambda}_y(x-\xi)$ is $\Delta_g$-harmonic on $\HH$.
\end{lemma}

\begin{proof}
Substitute
$X-t=(\lambda/\mu)(x-\xi)$ and
$dt=(\lambda/\mu)d\xi$ into
\eqref{eq:standard-poisson-kernel-X}.  This gives
\eqref{eq:poisson-density-transport} and
\eqref{eq:aeg-poisson-kernel-x}.  Positivity is immediate, and normalization
follows either from the same change of variables or from
\[
 \frac1\pi\int_{\R}\frac{h}{s^2+h^2}\,ds=1
 \qquad(h>0).
\]
Finally, $P_y(X-t)$ is Euclidean harmonic away from its boundary pole.
Equation~\eqref{eq:hyperbolic-laplace-beltrami} transports that assertion to
$\Delta_gP^{\mu,\lambda}_y=0$.
\end{proof}

\subsection{The ideal-boundary Dirichlet theorem}

We write $C_0(\R)$ for the continuous functions satisfying
$\varphi(\xi)\to0$ as $|\xi|\to\infty$.  Equivalently, such a function
extends continuously to $\R\cup\{\infty\}$ with value zero at $\infty$.

\begin{theorem}[$C_0$ ideal-boundary Dirichlet problem]
\label{thm:C0-poisson-dirichlet}
For every $\varphi\in C_0(\R)$, the function
\begin{equation}\label{eq:poisson-extension-x}
 (\mathcal P_{\mu,\lambda}\varphi)(x,y)
 :=\int_{\R}P^{\mu,\lambda}_y(x-\xi)\varphi(\xi)\,d\xi
\end{equation}
has the following properties.
\begin{enumerate}
\item It belongs to $C^\infty(\HH)\cap
 C(\overline\HH^{\,\mathrm{conf}})$ and is $\Delta_g$-harmonic.
\item It takes the boundary values
\[
 (\mathcal P_{\mu,\lambda}\varphi)(\xi,0)=\varphi(\xi)
 \quad(\xi\in\R),
 \qquad
 (\mathcal P_{\mu,\lambda}\varphi)(\infty)=0.
\]
\item It satisfies the contraction estimate
\begin{equation}\label{eq:poisson-sup-contraction}
 \|\mathcal P_{\mu,\lambda}\varphi\|_{L^\infty(\HH)}
 \leq\|\varphi\|_{L^\infty(\R)}.
\end{equation}
\item It is the unique function in
$C^2(\HH)\cap C(\overline\HH^{\,\mathrm{conf}})$ having these boundary
values and satisfying $\Delta_gU=0$.
\end{enumerate}
\end{theorem}

\begin{proof}
Put
\[
 \psi(t)=\varphi\left(\frac{\mu}{\lambda}t\right).
\]
Then $\psi\in C_0(\R)$, and Lemma~\ref{lem:poisson-kernel-transport}
rewrites \eqref{eq:poisson-extension-x} as
\begin{equation}\label{eq:poisson-extension-X}
 U(X,y)=\int_{\R}P_y(X-t)\psi(t)\,dt.
\end{equation}
For fixed $y>0$, derivatives of the kernel are integrable after localization,
and the usual differentiation argument shows that $U$ is smooth and
Euclidean harmonic.  Equivalently, one may first prove this for compactly
supported smooth data and use local convergence of the differentiated
kernels.  Equation~\eqref{eq:hyperbolic-laplace-beltrami} then gives
$\Delta_gU=0$.  Positivity and unit mass of the kernel prove
\eqref{eq:poisson-sup-contraction}.

The kernels $P_y$ form an approximate identity.  Hence, at each
$X_0\in\R$,
\[
 U(X,y)-\psi(X_0)
 =\int_{\R}P_y(X-t)\bigl(\psi(t)-\psi(X_0)\bigr)\,dt
 \longrightarrow0
\]
as $(X,y)\to(X_0,0)$.  To see the behavior at infinity without assuming
differentiability of the data, choose $\psi_0\in C_c(\R)$ with
$\|\psi-\psi_0\|_\infty<\varepsilon$.  The contraction estimate reduces the
claim to $\psi_0$.  If $\operatorname{supp}\psi_0\subset[-R,R]$, then its
Poisson integral tends to zero as $|X+iy|\to\infty$.  Indeed, for large $y$
it is bounded by $\|\psi_0\|_{L^1}/(\pi y)$; for large $|X|$, maximizing
$y/((X-t)^2+y^2)$ over $y>0$ gives a bound of order
$(|X|-R)^{-1}$.  Thus $U$ extends continuously to the compactification with
value zero at $\infty$.

For uniqueness, let $V$ be the difference of two solutions.  A Cayley map
sends $\HH$ to the unit disc and $\overline\HH^{\,\mathrm{conf}}$ to the
closed disc.  The pullback of $V$ is harmonic in the disc, continuous on its
closure, and zero on the entire boundary circle.  The maximum principle
applied to its real and imaginary parts gives $V=0$.
\end{proof}

\begin{corollary}[Continuous data with a value at infinity]
\label{cor:continuous-compactified-dirichlet}
Let $\varphi\in C(\R\cup\{\infty\})$, so that
\[
 L:=\varphi(\infty)=\lim_{|\xi|\to\infty}\varphi(\xi).
\]
Then
\[
 U=L+\mathcal P_{\mu,\lambda}(\varphi-L)
\]
is the unique harmonic function continuous on
$\overline\HH^{\,\mathrm{conf}}$ with boundary value $\varphi$.
\end{corollary}

\begin{proof}
The function $\varphi-L$ belongs to $C_0(\R)$, while constants are harmonic.
Apply Theorem~\ref{thm:C0-poisson-dirichlet}.
\end{proof}

\begin{warning}[Why the function class matters]
\label{warn:dirichlet-uniqueness-class}
If neither boundedness nor a condition at the ideal point is imposed,
uniqueness fails.  The function $U(x,y)=y$ is nonzero and harmonic, extends
continuously with value zero to every finite boundary point, and is unbounded
at the ideal point.  A separately prescribed bounded boundary value at the
single point $\infty$ also cannot alter a bounded Poisson solution: that point
has zero harmonic measure.  Continuous data on the compactified boundary must
have the same limit along both ends of the real line.
\end{warning}

The theorem is deliberately a boundary theorem, not an $L^2$ spectral
statement.  Unless its boundary data vanish in a sufficiently strong sense, a
Poisson extension is generally not in $L^2(\HH,d\vol_g)$ because the
hyperbolic volume density grows like $y^{-2}$ near the ideal boundary.
